\documentclass[11pt]{article}
\title{A guard on new lift payload writes does not retire older ones}
\date{26 November 2026}
\begin{document}
\maketitle

\paragraph{Modeled gate.}
The earlier source replay distinguishes advertised session from the
physical input to the simulator's per-payload velocity update.
Use that exact code in a serial test fixture, and put an authored
guard immediately before a proposed nonzero payload-write batch.
For this test the building operator can inspect the selected payloads,
enforce a permit check, and prevent the whole new batch. Each fleet's
scoped permit is bound to a modeled body, one trip instance, and a
closed range of effect cuts. The operator checks the permitted effect
of this instance, not a grant to any operation with the same subject.
The fixture has a known unique trip and immutable association to
the current input command for each run. A deployed binding between
trip identity and native command instance is not supplied.
Only the modeled velocity-command write is covered. Actual platform
motion could affect a physically present body even when it has
no velocity-command component.

At the guard and write cut, both reads of the selected shape and
positions use the same serial component state. Each recipient selected
by the native loop must have an active matching permit covering that
cut. Under this serialization, inclusion is both necessary and sufficient
for every new write in the batch to be authorized. The model refuses
the entire batch instead of suppressing writes to particular bodies
while claiming a shared ride still moves only allowed occupants.
If poses, shape, component identities, command binding or valid
permissions can change between the guard and use, this reasoning
does not apply without a binding or synchronization mechanism.

\paragraph{Replayed separation.}
An issue-time lookup allows the proposed trip while only the new
owner's modeled payload is present. A late unpermitted entry
blocks the batch when the effect-side check sees it. With trip-specific
consent an old payload stays onsite and the actual selected loop
writes to both payloads, without claiming the previous duty has
settled. A valid permit for a different trip does not count. The
effect-side guard also distinguishes the inclusive validity limits
of the correct trip's permit.

Next, revoke authorization after a successful write in the fixture.
The guard prevents another nonzero write. It does not clear the
earlier velocity values stored on either payload. Even passing
zero cabin velocity to the selected native branch skips its
per-payload write. The later values are unchanged in this serial
component fixture. This demonstrates why denying a new invocation
or observing a zero input to this branch does not prove an older
command retired. No conclusion follows about how long a full
Gazebo downstream system would retain or consume such values.
A revoke-and-settle interface needs independent knowledge of
actuator quiescence or a guaranteed actuator-side fence
covering already-issued writes and physical motion.

\paragraph{Limit.}
This is a guard on an executed source-level write selector, not
a lift safety controller, physical actuator test, authorization
system for fleet owners, or measurement of egress. An equally
informed direct effect check has the same result. No physical
clearance, service bound, independent held-out transfer, or
contract acceptance is claimed.
\end{document}
